\ResearchStyle
\PartDivider{PART III}{Related Work}{Compare MAC-only, TCP-only, and earlier cross-layer approaches.}

\begin{frame}{Related Work Map}
\begin{center}
\textbf{The paper organizes prior approaches as layered or cross-layer designs.}
\end{center}

\begin{columns}[T,onlytextwidth]
\column{0.48\textwidth}
\textbf{Layered design}
\begin{itemize}
\item MAC-focused studies: EIFS, backoff, and CW control.
\item TCP/network-focused studies: NRED and TCP window/rate control.
\end{itemize}
\column{0.48\textwidth}
\textbf{Cross-layer design}
\begin{itemize}
\item Exchange information across protocol layers.
\item Use MAC information to adapt TCP behavior.
\end{itemize}
\end{columns}

\vspace{0.25cm}
\KeyPoint{Prior work either solves problems independently at each layer or exchanges information across layers.}
\end{frame}

\begin{frame}{MAC Work: EIFS and Fair Backoff}
\begin{columns}[T,onlytextwidth]
\column{0.48\textwidth}
\textbf{Flexible EIFS (Li et al.)}
\begin{itemize}
\item Fixed EIFS can produce unfair allocation.
\item Flexible EIFS uses sensing-range frame information.
\item Spatial reuse can prevent reliable recognition of that information.
\end{itemize}
\column{0.48\textwidth}
\textbf{Fair backoff (Fang et al.)}
\begin{itemize}
\item Estimate own and neighboring channel utilization.
\item Adjust backoff toward statistical fair access.
\item The target fair-access value must be predefined and depends on topology.
\end{itemize}
\end{columns}
\end{frame}

\begin{frame}{MAC Work: CW Control}
\PaperClaim{Earlier CW-control studies seek per-flow fairness, but the CATRA paper highlights their single-hop focus.}
\begin{itemize}
\item Changing the contention window changes the expected access delay.
\item Single-hop assumptions simplify which flows and stations share the channel.
\item Multi-hop networks introduce forwarding, overlapping contention domains, and hidden carrier-sensing flows.
\end{itemize}
\end{frame}

\begin{frame}{TCP and Network-Layer Approaches}
\begin{columns}[T,onlytextwidth]
\column{0.48\textwidth}
\textbf{Neighborhood RED}
\begin{itemize}
\item Observe neighborhood channel usage.
\item Detect local congestion.
\item Drop packets according to flow usage to improve fairness.
\end{itemize}
\column{0.48\textwidth}
\textbf{TCP window limitation}
\begin{itemize}
\item Conventional TCP windows can become too large.
\item Excess packets reduce spatial reuse and increase contention.
\item Prior work argues that an optimal TCP rate/window may exist.
\end{itemize}
\end{columns}
\end{frame}

\begin{frame}{Earlier Cross-Layer TCP Methods}
\begin{itemize}
\item TCP-CL reports MAC link failure using the retry-limit state, but heavy contention can be mistaken for link failure.
\item TCP-FEW limits congestion-window growth to avoid unnecessary contention.
\item Channel-efficiency methods derive a better TCP rate from MAC information.
\item The CATRA paper argues that these approaches do not focus on wireless MAC contention or cannot exactly identify the relevant MAC channel efficiency at an intermediate node.
\end{itemize}
\vspace{0.35cm}
\KeyPoint{CATRA first determines a suitable CW at MAC, then applies TCP rate adaptation.}
\end{frame}

\begin{frame}{Related Work Comparison}
\small
\begin{tabular}{>{\raggedright\arraybackslash}p{2.1cm}>{\raggedright\arraybackslash}p{3.3cm}>{\raggedright\arraybackslash}p{5.1cm}}
\toprule
\textbf{Category} & \textbf{Control} & \textbf{Gap identified by CATRA} \\
\midrule
MAC / layered & EIFS, backoff, CW & Topology-dependent fair target; most CW studies are single-hop. \\
TCP / network & NRED, cwnd/rate & Treats TCP/network behavior, not MAC unfairness first. \\
Cross-layer & MAC information $\rightarrow$ TCP & No MAC-contention focus; intermediate-node efficiency is not exact. \\
CATRA & Suitable CW + TCP rate & Suitable CW per station, then TCP rate per flow. \\
\bottomrule
\end{tabular}
\end{frame}
